\section{Template Structure}

The class exposes a small set of metadata commands in addition to the standard \LaTeX{} \verb|\title|, \verb|\author|, and \verb|\date| commands.

\begin{lstlisting}[language=TeX,caption={Minimal metadata block.}]
\documentclass{nuspaper}
\title{Your Paper Title}
\shorttitle{Running Header Title}
\date{2026-06-15}
\arxivstamp{arXiv:XXXX.XXXXX [cs.LG] 15 Jun 2026}
\author[1,*]{First Author}
\affil[1]{National University of Singapore}
\paperlinks{\FrontLink{Project Page}{Project Website}{https://example.com}}
\titlefootnote{* Equal contribution.}
\end{lstlisting}

\subsection{Recommended repository layout}

Keep the manuscript modular but not over-engineered. A practical structure is:

\begin{lstlisting}[language=bash,caption={Suggested file layout.}]
nuspaper-template/
|-- main.tex
|-- nuspaper.cls
|-- refs.bib
|-- sections/
|   |-- 01_introduction.tex
|   |-- 02_method.tex
|   `-- 03_experiments.tex
|-- figures/
|-- tables/
`-- latexmkrc
\end{lstlisting}

\subsection{Mathematical environments}

Standard theorem-like environments are included. For example:

\begin{definition}[House-style invariant]
A paper template is \emph{venue-compatible} if its scientific content can be moved to an official submission template without rewriting the mathematical notation, citations, labels, figures, or tables.
\end{definition}

\begin{proposition}
A small class file plus standard manuscript body files is easier to maintain than a large, paper-specific preamble copied across projects.
\end{proposition}

\begin{proof}
The class file centralizes typography and layout decisions. The manuscript body stays mostly semantic, so later migration to a venue template only requires replacing the class and adapting a small number of metadata commands.
\end{proof}
